% ============================================================================
% Fichier  : partie4/P04_C19_rapport_forensique.tex
% Titre    : Rédaction du Rapport Forensique
% Auteur   : MINKA MI NGUIDJOÏ Thierry Emmanuel
% Version  : 1.0 -- Juin 2026
% Statut   : RÉDIGÉ
% Sources  : Ch. 8 (custody, admissibilité), Ch. 7 (NIST SP 800-86
%            distinction constatation/analyse), Ch. 12 (droit camerounais,
%            qualification), Ch. 13 (déontologie, témoignage expert)
%            Affaires fil rouge : MBONGO, MVOM, WOURI, ONAMBELE
% Positionnement : Chapitre-synthèse de la Partie IV. Toute l'investigation
%                  technique des chapitres précédents se traduit ici en
%                  document juridiquement opposable. C'est le livrable final
%                  de l'investigateur -- ce que le tribunal reçoit.
% ============================================================================

\chapter{Rédaction du Rapport Forensique}
\label{chap:rapport}

\epigraph{%
    « Un rapport forensique qui ne résiste pas à un contre-interrogatoire
    ne vaut pas mieux qu'une analyse technique jamais documentée. »%
}{--- Principe de la défendabilité forensique}

\niveauChapitre{Intermédiaire}{%
    Chapitres~\ref{chap:custody} (admissibilité) et
    \ref{chap:meth-std} (distinction constatation/analyse)
    indispensables. Ce chapitre synthétise toute la Partie~IV
    en document juridiquement opposable.%
}

\begin{boxobjectifs}
\begin{itemize}
    \item Maîtriser l'architecture en trois niveaux du rapport
          forensique~: synthèse exécutive, rapport technique,
          dossier probatoire.
    \item Appliquer rigoureusement la distinction entre
          \termecle{Constatations} (faits vérifiables) et
          \termecle{Analyse} (interprétations avec niveau de
          confiance).
    \item Structurer chaque section selon les exigences
          d'admissibilité du droit camerounais et les standards
          forensiques internationaux (NIST~SP~800-86).
    \item Rédiger une synthèse exécutive lisible par un magistrat
          sans formation technique.
    \item Anticiper les attaques en contre-interrogatoire et
          structurer les réponses dans le rapport lui-même.
    \item Produire un rapport complet sur le cas MVOM en appliquant
          l'intégralité des techniques présentées dans ce cours.
\end{itemize}
\end{boxobjectifs}

\bigskip

% ============================================================================
\section*{Pourquoi le rapport est le vrai livrable du cours}
% ============================================================================

Toute l'investigation technique des Chapitres~\ref{chap:for-sys}
à~\ref{chap:for-iot} ne produit qu'une chose visible pour le tribunal~:
le rapport forensique. Une analyse brillante mal documentée ne vaut
rien. Un rapport solide sur une analyse modeste peut convain­cre un juge.

Le rapport forensique est le seul contact direct entre votre travail
et la décision judiciaire. C'est aussi le document qui sera scruté
mot à mot par l'avocat de la défense. Ce chapitre traite donc la
rédaction non comme une formalité administrative mais comme
une compétence technique à part entière.

\separateur

% ============================================================================
\section{Architecture en Trois Niveaux}
\label{sec:architecture-rapport}
% ============================================================================

Conformément aux principes du PIU (Ch.~\ref{chap:PIU}, étape~E12) et
aux standards NIST~SP~800-86, le rapport forensique complet comprend
trois niveaux de profondeur correspondant à trois audiences distinctes.

\begin{tableaucours}{p{3.0cm}p{3.5cm}p{5.0cm}}{%
    \tableauHeader{Niveau} &
    \tableauHeader{Audience} &
    \tableauHeader{Contenu et longueur indicative}
}
\textbf{Niveau~1} --- Synthèse Exécutive &
    Magistrat, partie civile, direction &
    2--4~pages. Pas de technique. Conclusions, qualifications, preuves
    clés, recommandations. Compréhensible par un non-technicien. \\
\textbf{Niveau~2} --- Rapport Technique &
    Magistrat, avocat, expert adverse &
    20--100~pages. Méthodologie, constatations horodatées, analyse.
    Chaque affirmation référencée à une preuve. \\
\textbf{Niveau~3} --- Dossier Probatoire &
    Expert adverse, tribunal &
    Annexes~: images E01, logs bruts, captures d'écran, formulaires
    de custody. Référencés depuis le Niveau~2. \\
\end{tableaucours}
\vspace{-0.8em}
\begin{center}{\small\itshape Tableau~19.1 --- Architecture en trois niveaux du rapport forensique}\end{center}

\begin{boiteRemarque}{Rédiger pour le bon destinataire}
Le Niveau~1 est rédigé pour une personne qui ne sait pas ce qu'est
un hash SHA-256. Le Niveau~2 est rédigé pour un technicien
ou un avocat compétent qui \emph{voudra} comprendre la méthode.
Le Niveau~3 est rédigé pour un expert qui \emph{voudra vérifier}
chaque affirmation. Ces trois documents doivent être cohérents~:
la même conclusion doit apparaître dans les trois niveaux, à des
niveaux de détail différents.
\end{boiteRemarque}

% ============================================================================
\section{Règle Fondamentale~: Constatation vs Analyse}
\label{sec:constatation-analyse}
% ============================================================================

La distinction entre \textbf{Constatation} et \textbf{Analyse} est
la règle la plus importante de la rédaction forensique. Issue de
NIST~SP~800-86 (Ch.~\ref{chap:meth-std}), elle protège l'expert
des attaques méthodologiques les plus fréquentes.

\begin{boxdef}
\textbf{Constatation}~: énoncé factuel, objectif, vérifiable
indépendamment par tout technicien utilisant les mêmes outils
sur les mêmes données. Doit pouvoir être confirmé ou infirmé
sans ambiguïté. N'inclut aucune inférence.

\textbf{Analyse}~: interprétation des constatations dans leur
contexte. Comprend obligatoirement~: un niveau de confiance explicite
(élevé / modéré / faible) et les hypothèses alternatives envisagées
et réfutées (E9bis PIU).
\end{boxdef}

\begin{tableaucours}{p{5.5cm}p{5.5cm}}{%
    \tableauHeader{Constatation (correct)} &
    \tableauHeader{Mélange interdit (à éviter)}
}
« Le fichier \texttt{enc\_all.ps1} est présent dans
\texttt{\%TEMP\%} avec un horodatage de création du
13/03/2026~à~02h16:50 selon le \texttt{\$UsnJrnl}. » &
« Le fichier \texttt{enc\_all.ps1} a été placé par l'attaquant
dans \texttt{\%TEMP\%}. » (mélange fait + interprétation) \\
« Le journal Prefetch enregistre une exécution de
\texttt{enc\_all.ps1} à 02h17:44, avec un compteur
d'exécution de~1. » &
« L'attaquant a exécuté le ransomware à 02h17. »
(l'identité de l'attaquant n'est pas dans le Prefetch) \\
« Le \texttt{\$UsnJrnl} enregistre la modification de
847~fichiers \texttt{.docx} entre 02h18:31 et 02h19:02. » &
« L'attaquant a chiffré 847~fichiers en 31~secondes. »
(``chiffré'' est une interprétation~; ``modifié'' est le fait) \\
\end{tableaucours}
\vspace{-0.8em}
\begin{center}{\small\itshape Tableau~19.2 --- Exemples constatation vs mélange interdit}\end{center}

\begin{boxalert}
\textbf{La règle pratique~: le test du ``mais si''}

Avant d'écrire une phrase dans la section Constatations, posez-vous~:
``Mais si je me trompe sur l'interprétation, la phrase est-elle
encore vraie~?''

``\textit{Le fichier enc\_all.ps1 est présent dans \%TEMP\%}''
$\to$ vrai indépendamment de toute interprétation. \textbf{Constatation.}

``\textit{L'attaquant a déposé enc\_all.ps1 dans \%TEMP\%}''
$\to$ suppose que l'auteur du dépôt est identifié. Si l'identification
échoue, la phrase est fausse. \textbf{Analyse, pas constatation.}
\end{boxalert}

% ============================================================================
\section{Structure Détaillée du Rapport Technique (Niveau~2)}
\label{sec:structure-rapport}
% ============================================================================

\subsection{Page de garde et informations administratives}

\begin{lstlisting}[language=bash_forensique,
    caption={Métadonnées obligatoires --- page de garde}]
RAPPORT D'EXPERTISE FORENSIQUE NUMERIQUE
=========================================

Numero de dossier    : DOS-2026-034
Numero de rapport    : REP-DOS-2026-034-v1.0
Date de redaction    : 20 mars 2026
Affaire              : Opération MVOM -- ransomware TRANSBOIS SA

EXPERT :
Nom                  : [Nom Prénom]
Titre/Grade          : OPJ / Expert judiciaire agrée [si applicable]
Matricule            : PNC XXXX-X
Organisme            : DCPJ Douala
Contact              : [email officiel]  [telephone]

COMMETTANT :
Autorité mandante    : Parquet du TGI de Douala
Référence judiciaire : [Numéro de réquisition / ordonnance]
Date de la mission   : 13 mars 2026

EVIDENCE(S) ANALYSEE(S) :
EV-001 : Dell Inspiron 15 SN-GF8T3Y2 (image E01)
         Hash MD5 : a3f7b29c8e1d4f56a2c8b7e3d9f0a142
         Hash SHA-256 : 3d9c8f2b7e4a1c6d5f8e2b9a...
EV-001bis : Dump RAM 16 Go (MBONGO_RAM_20260315.mem)
         Hash SHA-256 : 4e8f3a1b2c9d7e5f6a0b1c2d...

VERSION : v1.0 -- rapport initial
          [v2.0 si rapport complémentaire requis]
\end{lstlisting}

\subsection{Section~1~: Contexte et Mission}

\begin{lstlisting}[language=bash_forensique,
    caption={Section Contexte et Mission --- exemple MVOM}]
1. CONTEXTE ET MISSION
======================

1.1 Contexte factuel
Le 13 mars 2026, TRANSBOIS SA (société d'exploitation forestière,
Douala) a signalé une attaque ransomware ayant paralysé son système
informatique. 847 fichiers de travail ont été rendus inaccessibles.

1.2 Mission
Conformément à la réquisition judiciaire N° [XXX] du Parquet du TGI
de Douala en date du 13 mars 2026, l'expert soussigné a été mandaté
pour :
  (a) Établir les faits techniques liés à l'incident du 13 mars 2026 ;
  (b) Identifier, si possible, l'origine de l'attaque ;
  (c) Évaluer l'étendue des dommages.

1.3 Périmètre et limitations
L'analyse porte exclusivement sur l'evidence EV-001 (image E01 du
serveur compromis) et EV-001bis (dump RAM).
Limitation : la capture réseau contemporaine à l'incident n'était
pas disponible. Les logs de pare-feu MTN pour la période ont été
requis séparément (réquisition N° [YYY] du 14 mars 2026 ; réponse
attendue).
\end{lstlisting}

\subsection{Section~2~: Méthodologie}

\begin{lstlisting}[language=bash_forensique,
    caption={Section Méthodologie --- exemple}]
2. MÉTHODOLOGIE
===============

2.1 Standards appliqués
- ISO/IEC 27037:2012 (collecte et préservation de l'evidence
  numérique)
- NIST SP 800-86 (intégration de la forensique dans la réponse
  aux incidents)
- RFC 3227 (ordre de collecte selon la volatilité)

2.2 Outils utilisés (versions exactes)
| Outil          | Version      | Usage                           |
| FTK Imager     | 4.7.3.81     | Acquisition de l'image E01      |
| Autopsy        | 4.21.0       | Analyse du système de fichiers  |
| Volatility 3   | 2.4.1        | Analyse mémoire RAM             |
| RegRipper      | 20230510     | Analyse du registre Windows     |
| PECmd          | 1.5.0.0      | Analyse des fichiers Prefetch   |
| MFTECmd        | 1.2.2.1      | Analyse du $MFT et $UsnJrnl     |
| log2timeline   | 20230724     | Construction de la timeline     |
| Systeme hôte   | Ubuntu 22.04 | Environnement d'analyse         |

2.3 Intégrité des données
Avant toute analyse, les hash SHA-256 des images E01 ont été
recalculés et comparés aux valeurs enregistrées lors de
l'acquisition (formulaire custody EV-001). Résultat :
CORRESPONDANCE CONFIRMÉE.
\end{lstlisting}

\subsection{Section~3~: Constatations}

La section Constatations est la plus importante du rapport~: elle
constitue la base de toute conclusion judiciaire. Elle ne contient
\textbf{que} des faits vérifiables, horodatés et référencés à une source.

\begin{lstlisting}[language=bash_forensique,
    caption={Section Constatations --- structure et exemple MVOM}]
3. CONSTATATIONS
================

[Structure recommandée : une constatation = un fait = une source]

3.1 État général du système
C.01 : Le système d'exploitation est Windows Server 2022 Standard,
       version 10.0.20348 Build 20348.
       [Source : Volatility 3, windows.info]

C.02 : Le compte utilisateur "admin" a été le dernier compte connecté
       sur ce serveur, avec une session active au moment de
       l'acquisition de la RAM.
       [Source : Volatility 3, windows.logonsessions]

3.2 Activité réseau identifiée
C.03 : Le journal journald enregistre 7 tentatives de connexion SSH
       échouées depuis l'IP 185.220.XX.XX entre 02h13:45 et 02h14:08.
       [Source : journald, extrait en Annexe A, lignes 1247-1253]

C.04 : Le journal journald enregistre une connexion SSH réussie au
       compte "root" depuis l'IP 185.220.XX.XX à 02h14:09.
       [Source : journald, extrait en Annexe A, ligne 1254]

3.3 Fichiers et activité système
C.05 : Le $UsnJrnl enregistre la création du fichier
       "enc_all.ps1" dans C:\Users\admin\AppData\Local\Temp\
       à 02h16:50.
       [Source : MFTECmd, fichier $UsnJrnl, Annexe B]

C.06 : Le journal Prefetch enregistre une exécution de "enc_all.ps1"
       à 02h17:44. Le compteur d'exécution est 1 (première exécution).
       Les fichiers accédés lors de cette exécution incluent
       "encrypted_files_list.txt".
       [Source : PECmd, ENC_ALL.PS1-3F8A2B91.pf, Annexe C]

C.07 : Le $UsnJrnl enregistre la suppression du fichier "enc_all.ps1"
       à 02h18:30.
       [Source : MFTECmd, fichier $UsnJrnl, Annexe B]

C.08 : Le $UsnJrnl enregistre la modification de 847 fichiers
       d'extension ".docx" entre 02h18:31 et 02h19:02. Ces
       fichiers portent désormais l'extension ".locked".
       [Source : MFTECmd, fichier $UsnJrnl, Annexe B]

C.09 : L'analyse Volatility (windows.netscan) révèle une connexion
       TCP en état ESTABLISHED entre 10.0.0.15:49212 et
       185.220.XX.XX:4444 au moment de la capture mémoire.
       [Source : Volatility 3, windows.netscan, Annexe D]
\end{lstlisting}

\subsection{Section~4~: Analyse}

L'analyse interprète les constatations~; elle doit~:
(a)~citer les constatations sur lesquelles elle s'appuie~;
(b)~exprimer un niveau de confiance~;
(c)~mentionner les hypothèses alternatives et pourquoi elles ont été
écartées (E9bis PIU).

\begin{lstlisting}[language=bash_forensique,
    caption={Section Analyse --- exemple MVOM}]
4. ANALYSE
==========

4.1 Chronologie de l'incident (niveau de confiance : ÉLEVÉ)

La convergence temporelle des sources C.03 à C.09 suggère, avec
un niveau de confiance élevé, le scénario suivant :

  [02h13:45 - 02h14:08] L'IP 185.220.XX.XX a procédé à plusieurs
  tentatives d'authentification SSH (C.03), constituant un schéma
  cohérent avec une attaque par force brute.

  [02h14:09] Une connexion SSH réussie a été établie au compte
  root depuis cette IP (C.04). La proximité temporelle (21 secondes
  après la dernière tentative échouée) est cohérente avec un accès
  obtenu par force brute.

  [02h16:50 - 02h17:44] Un script PowerShell nommé "enc_all.ps1"
  a été créé (C.05) puis exécuté (C.06) via la session SSH active.

  [02h18:31 - 02h19:02] 847 fichiers .docx ont été modifiés et
  renommés en .locked (C.08), schéma caractéristique d'un chiffrement
  par ransomware.

4.2 Hypothèse alternative écartée (falsification, E9bis PIU)

H_alt : L'activité observée pourrait être le résultat d'une
opération administrative légitime par l'administrateur système.

Réfutation : (1) L'IP source (185.220.XX.XX) n'appartient pas
au réseau interne de TRANSBOIS SA selon le schéma réseau fourni ;
(2) Aucune fenêtre de maintenance n'était planifiée le 13/03/2026
à 02h00 (attestation TRANSBOIS SA en Annexe E) ; (3) Le script
"enc_all.ps1" n'est pas présent dans les procédures documentées
de TRANSBOIS SA. Cette hypothèse est donc écartée avec un niveau
de confiance élevé.

4.3 Qualification technique

La combinaison accès SSH non autorisé + déploiement de script de
chiffrement + modification massive de fichiers constitue, sur le
plan technique, une intrusion informatique suivie d'une attaque
par ransomware. La qualification juridique de ces faits relève
de l'autorité judiciaire mandante.
\end{lstlisting}

\begin{boxalert}
\textbf{La qualification juridique n'appartient pas à l'expert}

La section~4 de l'exemple ci-dessus se termine par~:
``\textit{La qualification juridique de ces faits relève de
l'autorité judiciaire mandante.}''

Cette formule est obligatoire. L'expert dit ce qu'il a trouvé
et ce que cela signifie techniquement. Il ne dit pas ``c'est
une infraction à l'art.~65 de la Loi~2010/012''. Ce serait
outrepasser sa mission et fragiliser son témoignage.
\end{boxalert}

\subsection{Section~5~: Conclusions et Recommandations}

\begin{lstlisting}[language=bash_forensique,
    caption={Section Conclusions --- synthèse et recommandations}]
5. CONCLUSIONS ET RECOMMANDATIONS
==================================

5.1 Réponses aux questions posées par l'autorité mandante

Q1 : "Établir les faits techniques liés à l'incident."
R1 : Les faits techniques établis sont les suivants [résumé de
     l'analyse 4.1]. Ces faits sont établis avec un niveau de
     confiance élevé, sur la base de la convergence de six sources
     indépendantes (journald, $UsnJrnl, Prefetch, UserAssist,
     Volatility netscan, Volatility filescan).

Q2 : "Identifier l'origine de l'attaque."
R2 : L'IP source identifiée est 185.220.XX.XX. L'attribution de
     cette IP à un individu spécifique nécessite une réquisition
     aux autorités compétentes du pays d'enregistrement de cette IP.
     Cette démarche est hors du périmètre de la présente expertise.

Q3 : "Évaluer l'étendue des dommages."
R3 : 847 fichiers .docx ont été chiffrés et sont actuellement
     inaccessibles. [Évaluation complémentaire de la valeur économique
     des données hors périmètre de cette expertise.]

5.2 Recommandations de sécurité
[À inclure si demandé par le mandant -- ne pas inclure systématiquement
car peut être perçu comme sortant du périmètre judiciaire]

5.3 Réserves
Le présent rapport est établi sur la base des éléments disponibles
à la date de rédaction. Il est susceptible d'être complété si les
logs de pare-feu requis (réquisition N° [YYY]) sont communiqués.
\end{lstlisting}

% ============================================================================
\section{La Synthèse Exécutive (Niveau~1)~: Rédiger pour un Magistrat}
\label{sec:synthese-executive}
% ============================================================================

La synthèse exécutive est lue en premier et détermine si le rapport
sera compris. Elle ne contient aucun terme technique sans définition~;
elle ne référence pas de hash ou de commandes.

\begin{lstlisting}[language=bash_forensique,
    caption={Modèle de synthèse exécutive --- cas MVOM}]
SYNTHÈSE EXÉCUTIVE
==================
(À l'attention du Procureur de la République)

AFFAIRE : Opération MVOM -- Ransomware TRANSBOIS SA
DATE DE L'INCIDENT : 13 mars 2026, entre 02h14 et 02h19
EXPERT : [Nom], [Titre]

L'ESSENTIEL EN TROIS POINTS :

1. UN ACCÈS NON AUTORISÉ A ÉTÉ ÉTABLI
   Dans la nuit du 12 au 13 mars 2026, entre 02h14 et 02h19,
   une personne non autorisée a accédé au système informatique
   de TRANSBOIS SA depuis l'adresse Internet 185.220.XX.XX.
   Cet accès a été obtenu en devinant le mot de passe du
   compte administrateur après plusieurs tentatives.

2. 847 FICHIERS ONT ÉTÉ CHIFFRÉS ET RENDUS INACCESSIBLES
   L'intrus a utilisé un programme malveillant (ransomware)
   pour rendre 847 documents de travail inaccessibles.
   Ces documents portent désormais l'extension ".locked".

3. LES PREUVES SONT SOLIDES ET CONVERGENTES
   Six sources d'information indépendantes (journaux système,
   historique des fichiers, analyse de la mémoire vive)
   convergent vers les mêmes conclusions, avec un niveau de
   confiance élevé.

QUALIFICATION : La qualification juridique des faits relève
de l'autorité judiciaire. L'expert note que les faits techniques
établis correspondent à la description d'un accès informatique
non autorisé suivi d'une atteinte à l'intégrité des données.

PIÈCE CLEF : Le journal [LOGFILE] enregistre l'entrée de
l'intrus à 02h14:09 depuis l'adresse Internet 185.220.XX.XX.
Ce journal est annexé au présent rapport (Annexe A).

POUR ALLER PLUS LOIN : voir Rapport Technique (Niveau 2),
Section 3 (Constatations) et Section 4 (Analyse).
\end{lstlisting}

% ============================================================================
\section{Annexes~: le Dossier Probatoire (Niveau~3)}
\label{sec:dossier-probatoire}
% ============================================================================

Les annexes constituent la preuve que le rapport est vérifiable.
Elles doivent être organisées, référencées et intègres.

\begin{tableaucours}{p{2.5cm}p{4.0cm}p{5.0cm}}{%
    \tableauHeader{Annexe} &
    \tableauHeader{Contenu} &
    \tableauHeader{Format et exigences}
}
\textbf{Annexe~A} & Extraits de logs cités dans les Constatations &
    Texte brut~; horodatages lisibles~; lignes surlignées si nécessaire \\
\textbf{Annexe~B} & Export CSV du \texttt{\$UsnJrnl} filtré &
    CSV~; colonnes~: Timestamp, Filename, Reason~; trié par date \\
\textbf{Annexe~C} & Rapport PECmd (Prefetch) &
    Export CSV ou texte de PECmd \\
\textbf{Annexe~D} & Captures d'écran Volatility &
    PNG horodaté~; commande visible~; résultat complet \\
\textbf{Annexe~E} & Attestation de l'entreprise victime &
    Document signé par le responsable TRANSBOIS~SA \\
\textbf{Annexe~F} & Formulaire de custody EV-001 et EV-001bis &
    Formulaires remplis et signés (Ch.~\ref{chap:custody}) \\
\textbf{Annexe~G} & Hash SHA-256 de toutes les pièces &
    Fichier texte~: SHA-256~<espace>~<nom du fichier>~; signé par l'expert \\
\end{tableaucours}
\vspace{-0.8em}
\begin{center}{\small\itshape Tableau~19.3 --- Organisation des annexes du dossier probatoire}\end{center}

\begin{boiteRemarque}{L'Annexe G~: le hash de toutes les pièces}
L'Annexe~G est la plus importante~: elle est le méta-garant de l'intégrité
de toutes les autres. Elle contient un hash SHA-256 de chaque fichier
annexé au rapport. Si un avocat affirme qu'une pièce a été modifiée,
l'Annexe~G permet de le réfuter immédiatement en recalculant le hash.
Produire l'Annexe~G signée par l'expert au début de l'audience.
\end{boiteRemarque}

% ============================================================================
\section{Pièges de Rédaction et Comment les Éviter}
\label{sec:pieges-redaction}
% ============================================================================

\begin{tableaucours}{p{3.8cm}p{3.5cm}p{4.5cm}}{%
    \tableauHeader{Piège} &
    \tableauHeader{Exemple fautif} &
    \tableauHeader{Formulation correcte}
}
\textbf{Certitude excessive} &
    ``L'attaquant est X.'' &
    ``L'IP source est 185.220.XX.XX. L'attribution à une personne
    physique nécessite des investigations complémentaires.'' \\
\textbf{Jargon non défini} &
    ``Le ransomware a modifié le MFT et les entrées USN.'' &
    ``Les métadonnées de fichier (registre chronologique du système)
    indiquent~\ldots'' puis définir les termes en note de bas de page. \\
\textbf{Conclusion sans preuve} &
    ``L'accès a été obtenu par force brute.'' &
    ``7~tentatives échouées (C.03) suivies d'un succès (C.04)
    constituent un schéma cohérent avec une attaque par force brute~;
    d'autres explications sont théoriquement possibles mais
    n'ont pas été corroborées.'' \\
\textbf{Source non citée} &
    ``Les logs montrent~\ldots'' &
    ``Le journal \texttt{journald} (Annexe~A, lignes 1247--1254)
    enregistre~\ldots'' \\
\textbf{Hypothèse non réfutée} &
    ``Le suspect est coupable.'' &
    Présenter et réfuter explicitement les hypothèses alternatives
    (cf. E9bis PIU). \\
\textbf{Qualif. juridique par l'expert} &
    ``Ce fait constitue une infraction à l'art.~65.'' &
    ``La qualification juridique relève de l'autorité mandante.
    L'expert note que~\ldots'' \\
\end{tableaucours}
\vspace{-0.8em}
\begin{center}{\small\itshape Tableau~19.4 --- Pièges de rédaction et formulations correctes}\end{center}

% ============================================================================
\section{Cas ONAMBELE~: le Rapport Sous Pression}
\label{sec:rapport-pression}
% ============================================================================

\begin{boxscenario}{Affaire ONAMBELE~: quand le rapport est contesté}
L'expert NKENG a remis son rapport sur l'affaire ONAMBELE (fraude
sur les marchés publics). La partie adverse produit une contre-expertise
affirmant que les logs analysés par NKENG pourraient avoir été
manipulés par TRANSBOIS SA \emph{après} l'incident.

L'avocat de la défense pose trois questions en audience~:

\textbf{Q1~:} ``Pouvez-vous prouver que les logs que vous avez
analysés n'ont pas été modifiés par votre client après l'incident~?''

\textbf{R (correcte)~:} ``Les logs ont été exportés depuis l'image
E01 (EV-001) dont les hash SHA-256 ont été calculés lors de
l'acquisition le 13 mars 2026 à 15h03 et sont documentés dans
le formulaire de custody (Annexe~F). La valeur de hash calculée
au moment de l'analyse le 15 mars 2026 est identique (Annexe~G).
Toute modification des logs après l'acquisition aurait changé
le hash et serait détectable. »

\textbf{Q2~:} ``Vous utilisez Volatility 3. N'est-il pas vrai que
ce logiciel modifie lui-même la mémoire pendant la capture~?''

\textbf{R (correcte)~:} « La capture mémoire (EV-001bis) a été
réalisée \emph{avant} toute autre opération sur le système, conformément
à RFC~3227 (ordre de volatilité~: RAM en premier). La modification
inhérente à l'outil de capture est documentée dans le formulaire
de custody sous la note~: \textit{``légère modification de la RAM
inhérente au processus de capture, conformément à ACPO Principe~2''}
(Annexe~F, ligne~12). Cette modification affecte quelques kilo-octets
et n'altère pas les artefacts analysés. »

\textbf{Q3~:} ``Votre rapport dit que l'attaque vient de l'IP
185.220.XX.XX. Comment savez-vous qu'il ne s'agit pas d'un VPN
utilisé par quelqu'un d'autre~?''

\textbf{R (correcte)~:} « Je ne l'affirme pas dans mon rapport.
La page~3 indique explicitement que l'attribution de cette IP à une
personne physique est hors périmètre de cette expertise et nécessite
une réquisition aux autorités compétentes. Je constate un fait~:
l'IP source des connexions SSH. L'identification du titulaire relève
d'une démarche distincte. »
\end{boxscenario}

\separateur

\begin{boxresume}
\textbf{Ce qu'il faut retenir de ce chapitre~:}
\begin{itemize}
    \item \textbf{Trois niveaux}~: Synthèse Exécutive (2--4 pages,
          non-technicien)~; Rapport Technique (20--100 pages, technicien
          ou avocat)~; Dossier Probatoire (annexes vérifiables).
          Cohérents entre eux~; rédigés pour des audiences différentes.

    \item \textbf{La règle d'or}~: Constatation = fait vérifiable
          (horodaté, sourcé). Analyse = interprétation (niveau de
          confiance + hypothèses alternatives réfutées). Ne jamais
          mélanger. Tester chaque phrase~: ``mais si je me trompe,
          la phrase est-elle encore vraie~?''

    \item \textbf{Chaque constatation cite sa source}~: ``[Source~:
          MFTECmd, fichier \$UsnJrnl, Annexe~B]''. Toujours. Sans
          exception.

    \item \textbf{Annexe~G obligatoire}~: hash SHA-256 de toutes les
          pièces annexées. C'est la preuve de l'intégrité du dossier
          probatoire lui-même.

    \item \textbf{La qualification juridique n'appartient pas à l'expert.}
          Finir la section Analyse par~: ``La qualification juridique
          de ces faits relève de l'autorité judiciaire mandante.''

    \item \textbf{Hypothèses alternatives obligatoires}~: présenter et
          réfuter au moins une hypothèse alternative par conclusion
          majeure. Sans cette réfutation, la conclusion est une
          affirmation, pas une analyse.

    \item \textbf{Synthèse exécutive}~: zéro jargon non défini~;
          trois points maximum~; une pièce clé identifiée~;
          une référence vers le rapport technique pour approfondir.
\end{itemize}
\end{boxresume}

\bigskip

\begin{boxexo}[title={\ding{45}\quad Exercice~19.1 --- Identifier constatation vs analyse \niveaubase}]
Pour chacun des énoncés suivants, indiquez s'il s'agit d'une
Constatation, d'une Analyse, ou d'un mélange interdit~; reformulez
les mélanges~:
\begin{enumerate}[(a)]
    \item « Le fichier \texttt{malware.exe} a été exécuté par
          MBONGO le 15 janvier 2026 à 14h28. »
    \item « Le Prefetch \texttt{MALWARE.EXE-A1B2C3D4.pf} enregistre
          une exécution à 14h28:33, compteur~= 1. »
    \item « L'attaquant a tenté de couvrir ses traces en supprimant
          le fichier enc\_all.ps1. »
    \item « Le \texttt{\$UsnJrnl} enregistre la suppression du
          fichier enc\_all.ps1 à 02h18:30, soit 46~secondes après
          l'entrée \texttt{DataExtend|Close} (02h17:44). »
    \item « Les logs de connexion montrent un accès suspect
          depuis une IP russe. »
\end{enumerate}
\end{boxexo}

\begin{boxexo}[title={\ding{45}\quad Exercice~19.2 --- Rédiger une section Constatations \niveaumoyen}]
À partir des données suivantes (extraites des analyses MVOM),
rédigez 5~constatations correctement structurées (fait + source
+ annexe)~:
\begin{itemize}
    \item UserAssist~: \texttt{AnyDesk.exe} exécuté 3~fois,
          dernière à 02h10:15~;
    \item Volatility netscan~: connexion ESTABLISHED vers
          185.220.XX.XX:4444 à partir du PID~4892~;
    \item Volatility pslist~: PID~4892 = svchost.exe,
          PID parent~= 4888 (cmd.exe)~;
    \item \$UsnJrnl~: création de \texttt{encrypted\_files\_list.txt}
          à 02h18:31~;
    \item journald~: SSH Accepted password pour root depuis
          185.220.XX.XX à 02h14:09.
\end{itemize}
\end{boxexo}

\begin{boxexo}[title={\ding{45}\quad Exercice~19.3 --- Rapport complet \niveauexpert}]
Produisez le rapport forensique complet sur l'Opération MVOM
(Chapitre~\ref{chap:gr-affaires}) en appliquant toutes les
techniques de ce cours~:
\begin{enumerate}[(a)]
    \item \textbf{Synthèse Exécutive}~: 1~page maximum,
          non-technicien, trois points, une pièce clé.
    \item \textbf{Section Constatations}~: au moins 10~constatations
          couvrant~: l'accès SSH, l'exécution du ransomware,
          le chiffrement des fichiers, et la tentative de suppression
          des traces. Chaque constatation cite sa source.
    \item \textbf{Section Analyse}~: chronologie de l'incident avec
          niveau de confiance explicite~; au moins une hypothèse
          alternative présentée et réfutée (E9bis PIU).
    \item \textbf{Section Conclusions}~: réponses aux 3~questions
          du mandant~; qualification technique sans qualification
          juridique.
    \item \textbf{Plan des Annexes}~: liste les pièces à annexer
          avec leur description et leur hash attendu.
\end{enumerate}
\textit{Note~: cet exercice peut être évalué comme devoir maison.}
\end{boxexo}

\bigskip

\begin{boxinfo}
\textbf{Transition.} La Partie~IV est terminée. Le cours dispose
maintenant de toutes les briques techniques~: investigation système,
réseau, mobile, cloud et IoT, plus la synthèse documentaire.

La \textbf{Partie~V} traite les contre-mesures~: Anti-Forensique
(techniques d'obfuscation et d'effacement~; comment les détecter),
puis les sujets avancés (OSINT, investigations spécialisées,
perspectives futures). Ces chapitres s'appuient sur la maîtrise
des techniques forensiques de la Partie~IV~--- la connaissance
de l'attaque est indissociable de la maîtrise de la défense.
\end{boxinfo}
